\section{结论、局限性与未来工作}

决赛运行时将专家预测视为可撤回、受安全和压力约束的建议：错误预取只对经过边界验证的内部页回收；驻留仅在压力、准确性和预算允许时准入；模型所有权与租约定义了后台工作和映射地址的共同生命周期。边界条件、重复认领、压力状态、指标解析和打包变体均有自动化门禁；正式实验进一步把每个成功运行绑定到源码、patched tree、模型、镜像与 cgroup 证据。

其局限性同样明确。页建议是 best-effort，内核和文件系统行为可能使性能波动；保守的页对齐会损失覆盖率；压力阈值和预测信号可能不适用于所有模型或工作负载；单一 2\,GiB/no-swap、短生成工作负载不能推广到其他设备或长上下文；所有配置都触及同一 2\,GiB 峰值，无法证明峰值内存下降；combined cold 回退 \FinalsCombinedColdRegressionPercent\%，说明组合并非天然协同。未来工作应扩展到多个受控内存档位和更长生成，分别测量真实 reclaim 命中、页错误与 I/O，再研究不会破坏地址与生命周期不变量的细粒度专家布局。图\ref{fig:reclaim-loop}仍是设计闭环，性能解释只接受新的受控证据。
